
When reinforcement learning agents learn to generate code, how much of each generated token actually matters? In a typical code solution, only a fraction of tokens---those that execute during test evaluation---directly influence the reward signal. The remaining tokens, including unused imports, unreached branches, and dead code, contribute nothing to the outcome yet receive identical gradient updates under standard policy gradient methods.

This observation points to a fundamental limitation in reinforcement learning from execution feedback for code generation. The sparse reward problem---where a binary test pass/fail signal provides no gradient for partial progress---has motivated numerous approaches to provide denser supervision. However, existing solutions often conflate two orthogonal factors: \textit{what} feedback is provided (compilation errors, test results, execution traces) and \textit{how} credit is assigned to individual tokens (episode-level rewards versus token-level masking).

Fine-Grained Optimization (FGO), introduced by StepCoder, addresses the credit assignment problem by masking non-executed code segments from gradient updates. The intuition is compelling: if a token never executes, it cannot have caused the test to pass or fail, so excluding it from learning should concentrate gradients on causally relevant code. Yet StepCoder introduces FGO alongside a curriculum learning strategy (CCCS), making it impossible to isolate which component drives the reported improvements.

This paper presents a controlled mechanism validation study of FGO for code generation RL. Rather than proposing a new method, the FGO mechanism is decomposed into three testable components and each is verified with explicit falsification criteria:

\begin{enumerate}
\item \textbf{Trace Collection (H-M1)}: Verification that Python's \texttt{sys.settrace} reliably captures execution traces during test evaluation.
\end{enumerate}

\begin{enumerate}
\item \textbf{Gradient Exclusion (H-M2)}: Verification that FGO's token masking correctly excludes non-executed tokens from gradient computation.
\end{enumerate}

\begin{enumerate}
\item \textbf{Credit Assignment (H-M3)}: Verification that concentrating gradients on executed tokens improves final performance.
\end{enumerate}

This decomposition enables precise failure localization. If trace collection fails, the problem lies in the execution monitoring infrastructure. If gradient exclusion fails, the masking implementation is incorrect. If credit assignment fails despite correct masking, the theoretical premise of FGO---that execution-aligned gradients improve learning---would be falsified.

The following contributions are made:

\begin{itemize}
\item \textbf{Mechanism Decomposition}: FGO is decomposed into three independently testable hypotheses (trace $\rightarrow$ mask $\rightarrow$ exclusion $\rightarrow$ credit), enabling controlled validation of each component.
\end{itemize}

\begin{itemize}
\item \textbf{Verification Protocol}: Explicit gate criteria with falsification thresholds are established, moving beyond aggregate performance metrics to mechanism-level validation.
\end{itemize}

\begin{itemize}
\item \textbf{Empirical Validation}: All three mechanism components are verified on HumanEval and MBPP benchmarks, demonstrating 100\% trace capture, verified gradient exclusion, and 10\% performance improvement in simulation.
\end{itemize}

\begin{itemize}
\item \textbf{Reusable Components}: Validated implementations of trace collection, token classification, and masked PPO loss are provided for future code RL research.
\end{itemize}
